\section{An exact comparison of equality, volume contraction, and its discarded square}

This calculation extends the positive Gaussian calibration of E35--E36
through two further source degrees. Every measure, phase, coefficient and
quotient basis is specified. It is a calibration of the same constructions;
the roots below are declared coefficients, with no assertion that they are
zeros of the arithmetic function.

Fix $\mu>0$, $\sigma>0$, $k=1$ and $c=1/2$. Put
\[
 d\nu(u)=\frac{\mu}{\sqrt{2\pi}\sigma}
             e^{-u^2/(2\sigma^2)}\,du,\qquad
 x=S-c,\qquad \chi(S)=x^2-\sigma^2,
 \qquad C=\mathbb C[S]/(\chi).
 \tag{TC1}
\]
The original source pairing is
$\langle P,Q\rangle=\int\overline{P(c+iu)}Q(c+iu)d\nu(u)$.
We use both the original quotient basis $(1,S)$ and the displayed basis
$(1,x)$. The exact coordinate matrix, sending columns in the latter basis
to columns in the former, is
\[
 L_{\mathrm{coord}}=\begin{pmatrix}1&-c\\0&1\end{pmatrix},\qquad
 A_S=\begin{pmatrix}0&\sigma^2-c^2\\1&2c\end{pmatrix},\qquad
 A_x=\begin{pmatrix}c&\sigma^2\\1&c\end{pmatrix},\qquad
 A_SL_{\mathrm{coord}}=L_{\mathrm{coord}}A_x.
 \tag{TC2}
\]
Here $A_x$ represents multiplication by $S$ in the basis $(1,x)$;
multiplication by $x$ is $A_x-cI$.
Thus $G_{N,x}=L_{\mathrm{coord}}^*G_{N,S}L_{\mathrm{coord}}$ and
$K_{N,S}=L_{\mathrm{coord}}K_{N,x}L_{\mathrm{coord}}^*$.
In particular $\det L_{\mathrm{coord}}=1$ preserves the
specified determinant, while the matrices and source polynomials retain
their full coordinate transformations.

Integration by parts, with the Gaussian boundary term zero, gives
$\int u^{2j}d\nu=\mu(2j-1)!!\sigma^{2j}$ and all odd moments zero:
indeed the zeroth integral is $\mu$, and differentiating
$u^{2j-1}e^{-u^2/(2\sigma^2)}$ gives the recurrence
$m_{2j}=(2j-1)\sigma^2m_{2j-2}$. The first four monic source
polynomials and their full squared norms are consequently
\[
 \begin{array}{c|c|c}
 j&p_j(S)&\omega_j\\\hline
 0&1&\mu\\
 1&x&\mu\sigma^2\\
 2&x^2+\sigma^2&2\mu\sigma^4\\
 3&x^3+3\sigma^2x&6\mu\sigma^6
 \end{array}.
 \tag{TC3}
\]
For example $p_2(c+iu)=-(u^2-\sigma^2)$ and
$p_3(c+iu)=-i(u^3-3\sigma^2u)$. Substitution of the moments
proves every pairwise orthogonality and every displayed norm; the phases
$-1$ and $-i$ are retained before taking the pairings.

In the basis $(1,x)$ the quotient columns $b_j=[p_j]_\chi$ are
\[
 b_0=\binom10,\quad b_1=\binom01,\quad
 b_2=\binom{2\sigma^2}0,\quad b_3=\binom0{4\sigma^2}.
 \tag{TC4}
\]
For completeness, the minimum-source metric is obtained by writing a
source polynomial as $\sum_{j=0}^N z_jp_j$. Its norm is
$\sum\omega_j|z_j|^2$ and its quotient is $\sum b_jz_j$.
The constrained minimizer is
$z_j=\omega_j^{-1}b_j^*G_{N,x}v$, where
$K_{N,x}=\sum b_jb_j^*/\omega_j$ and $G_{N,x}=K_{N,x}^{-1}$.
Indeed these coefficients have quotient $v$, and their inner product
with every quotient-kernel coefficient vector is zero. Pythagoras then
proves uniqueness and minimum norm $v^*G_{N,x}v$.
Set $V_N=\det G_{N,x}=\det G_{N,S}$. Directly from (TC3)--(TC4),
\[
 \begin{split}
 G_{1,x}&=\operatorname{diag}(\mu,\mu\sigma^2),\\
 G_{2,x}&=\operatorname{diag}(\mu/3,\mu\sigma^2),\\
 G_{3,x}&=\operatorname{diag}(\mu/3,3\mu\sigma^2/11),
 \end{split}\qquad
 (V_1,V_2,V_3)=\mu^2\sigma^2(1,1/3,1/11).
 \tag{TC5}
\]

This is also exactly the source/relation determinant formula. On the
original vertical line $\chi(c+iu)=-(u^2+\sigma^2)$, so its zeroth
relation norm and its next monic relation norm are
\[
 \nu_0=\int(u^2+\sigma^2)^2d\nu=6\mu\sigma^4,\qquad
 \nu_1=\int u^2(u^2+\sigma^2)^2d\nu=22\mu\sigma^6.
 \tag{TC6}
\]
The two relation columns $\chi,\chi x$ are orthogonal by parity.
For $n=2,3,4$, the source Gram determinant is
$\mathfrak D_n=\prod_{j<n}\omega_j$; the corresponding relation
determinants are $\mathfrak B_0=1$, $\mathfrak B_1=\nu_0$,
$\mathfrak B_2=\nu_0\nu_1$. These identities follow by the triangular
monic changes of basis. They give
$\mathfrak D_{N+1}/\mathfrak B_{N-1}=V_N$ for each $N=1,2,3$,
including the literal empty determinant at $N=1$.

Let $H_N=A_x+G_{N,x}^{-1}A_x^*G_{N,x}-2cI$. Multiplication of the
displayed matrices yields
\[
 H_1=\begin{pmatrix}0&2\sigma^2\\2&0\end{pmatrix},\qquad
 H_2=\begin{pmatrix}0&4\sigma^2\\4/3&0\end{pmatrix},\qquad
 \epsilon_1^2=4\sigma^2,\quad\epsilon_2^2=16\sigma^2/3.
 \tag{TC7}
\]
Each $H_N$ is self-adjoint in its stated positive metric; its squared
matrix is the corresponding scalar multiple of $I$. Thus these are
the squares of its operator norm, or equivalently
$\operatorname{Tr}(H_N^2)/2$. The spectral aggregate for $A_x$ is
$L_{\mathrm{spec}}=2\sigma$, since its eigenvalues are $c+\sigma,c-\sigma$.
Hence equality $\epsilon_1=L_{\mathrm{spec}}$ holds at the first admitted degree,
whereas $\epsilon_2>L_{\mathrm{spec}}$ holds at the next. No freedom to choose another
metric was used in either calculation.

To compute the auxiliary derivative at the original observation, replace
$d\nu$ by $e^{\theta u}d\nu$ for real $\theta$ and set
$\ell_N(\theta)=\log\det G_N(\theta)$ in either specified quotient basis.
The precise real-coordinate moment Grams are
$H_{rs}(\theta)=\int u^{r+s}e^{\theta u}d\nu$ and
$H^\chi_{rs}(\theta)=\int u^{r+s}(u^2+\sigma^2)^2e^{\theta u}d\nu$.
Their matrices at $-\theta$ are conjugated by
$\operatorname{diag}(1,-1,1,\ldots)$ from those at $\theta$.
The exact map to original monomials is
$S^j=\sum_{r=0}^j\binom jr c^{j-r}i^ru^r$.
On dimension $n$ its triangular determinant is $i^{n(n-1)/2}$,
of modulus one. It therefore identifies these determinants with the
original source and relation determinants, including the relation
multiplier $\chi$. Both determinants are even functions of $\theta$.
Their positive determinant ratio is therefore even, giving
$\ell_N'(0)=0$. Also
$\sigma_T:=\operatorname{Tr}((A_x-cI)/i)=0$.
This proves the vanishing of the phase subtraction at zero tilt; it
does not impose evenness at a nonzero tilt.

Finally set $\delta_N=V_N/V_{N-1}$ at the two degrees where both
metrics exist, and set $\mathcal R_2=V_1/V_3=(\delta_2\delta_3)^{-1}$.
Equations (TC3) and (TC5) give
\[
 \delta_2=1/3,\qquad \delta_3=3/11,\qquad
 a_3=\omega_3/\omega_2=3\sigma^2,\qquad \mathcal R_2=11.
 \tag{TC8}
\]
The complete Toda expression and its finite upper bound are respectively
\[
 a_3(1-\delta_2)(\delta_3^{-1}-1)=16\sigma^2/3,
 \qquad a_3\frac{(\mathcal R_2-1)^2}{4\mathcal R_2}
       =75\sigma^2/11.
 \tag{TC9}
\]
To identify the exact lost term, retain
$s=(\log3+\log(11/3))/2=\log(11)/2$ and
$t=(\log(11/3)-\log3)/2$. Then
$e^t=\sqrt{11}/3$ and $\cosh s=6/\sqrt{11}$, whence
\[
 a_3(e^t-\cosh s)^2=49\sigma^2/33
       =75\sigma^2/11-16\sigma^2/3.
 \tag{TC10}
\]
Thus the entire gap between the exact arithmetic-construction radius in
this specified calibration and its Toda upper expression is precisely
the retained imbalance square. The phase square is zero for the proved
parity reason. Both metric degrees, the quotient matrix, the source and
relation volumes, and the original mass $\mu$ occur explicitly in this
comparison.

There is also an explicit comparison with the spectral and orthogonal
projectors in the exterior equality proof. In the $(1,x)$ coordinates
let $v_+=(\sigma,1)^{\mathsf T}$ and $v_-=(-\sigma,1)^{\mathsf T}$.
They are eigenvectors of $A_x$ for $c+\sigma$ and $c-\sigma$. Directly,
\[
 v_+^*G_{1,x}v_-=0,\qquad
 v_+^*G_{2,x}v_-=2\mu\sigma^2/3.
 \tag{TC11}
\]
At $N=2$ the orthogonal projector to $\mathbb Cv_+$ and the algebraic
spectral projector along $\mathbb Cv_-$ are respectively
\[
 P=\frac{v_+v_+^*G_{2,x}}{v_+^*G_{2,x}v_+}
   =\begin{pmatrix}1/4&3\sigma/4\\1/(4\sigma)&3/4\end{pmatrix},
 \qquad
 Q=\frac{A_x-(c-\sigma)I}{2\sigma}
   =\begin{pmatrix}1/2&\sigma/2\\1/(2\sigma)&1/2\end{pmatrix}.
 \tag{TC12}
\]
These formulas prove $P^2=P$, $P^{\dagger_{G_{2,x}}}=P$, $Q^2=Q$,
and $\operatorname{im}P=\operatorname{im}Q=\mathbb Cv_+$.
Their exact kernels are
$\ker P=\mathbb C(-3\sigma,1)^{\mathsf T}=v_+^{\perp_{G_{2,x}}}$
and $\ker Q=\mathbb Cv_-$: direct multiplication proves the two
zero values, while $Pv_+=Qv_+=v_+$ proves rank one.
In particular $Pv_-=v_+/2$.
They also retain the change to original coordinates: both matrices are
sent to $L_{\mathrm{coord}}(\cdot)L_{\mathrm{coord}}^{-1}$ and paired
with $G_{2,S}=L_{\mathrm{coord}}^{-*}G_{2,x}L_{\mathrm{coord}}^{-1}$.
For $X=Q-P$, the exact Hilbert--Schmidt pairing gives
\[
 \|X\|_{\mathrm{HS},G_{2,x}}^2
 =\operatorname{Tr}(G_{2,x}^{-1}X^*G_{2,x}X)=1/3,
 \qquad
 \epsilon_2^2-L_{\mathrm{spec}}^2
      =L_{\mathrm{spec}}^2\|X\|_{\mathrm{HS},G_{2,x}}^2
      =4\sigma^2/3.
 \tag{TC13}
\]
For example the two diagonal entries of
$G_{2,x}^{-1}X^*G_{2,x}X$ are $1/4$ and $1/12$, which proves the
trace value without changing its pairing. The last equality follows
also from (TC7). Thus the two explicit gaps have different exact source
maps: (TC13) measures the spectral/orthogonal projector discrepancy,
while (TC10) measures the square discarded by the Toda upper bound.
Both maps and both positive quantities are calculated in the original
specified metric.
